\chapter*{Vorwort}
\addcontentsline{toc}{chapter}{Vorwort}

Dieses Skript ist als Lernskript zur Lehrveranstaltung \emph{Maschinelles Lernen}
gedacht. Es orientiert sich an der Themenliste der Vorlesung und des Proseminars und erklärt die Themen
von Grund auf.

Das Ziel ist nicht eine kurze Zusammenfassung, sondern ein verständlicher Aufbau der
wichtigsten Begriffe, Formeln, Herleitungen, Beispiele und Prüfungszusammenhänge.

\section*{Sprache und Begriffe}

Das Skript ist vollständig auf Deutsch geschrieben. Wichtige englische Fachbegriffe werden beim ersten Auftreten
in eckigen Klammern ergänzt, zum Beispiel \emph{überwachtes Lernen [supervised learning]},
\emph{Wahrscheinlichkeitsdichte [probability density function]} oder
\emph{Maximum-Likelihood-Schätzung [maximum likelihood estimation]}.

\section*{Aufbau der Kapitel}

Jedes Kapitel ist nach einem ähnlichen Schema aufgebaut:
\begin{itemize}
    \item Definitionen wichtiger Begriffe,
    \item intuitive Erklärung der Grundideen,
    \item zentrale Formeln,
    \item Herleitungen,
    \item Beispiele,
    \item prüfungsrelevante Hinweise.
\end{itemize}

\section*{Boxen}

In diesem Skript werden verschiedene Boxen verwendet:
\begin{itemize}
    \item \texttt{definitionbox}: wichtige Begriffe,
    \item \texttt{conceptbox}: zentrale Idee,
    \item \texttt{formulabox}: zentrale Formeln,
    \item \texttt{derivationbox}: Herleitungen,
    \item \texttt{examplebox}: Beispiele,
    \item \texttt{examtipbox}: Hinweise für die Prüfung.
\end{itemize}

\section*{Hinweis}

Das Skript ist als Arbeitsdokument gedacht. Es ersetzt nicht die Vorlesung,
sondern soll helfen, die Themen systematisch zu verstehen und zu wiederholen.
